\subsection{Computation, training, and inference}

Token IDs are mapped to learned vectors and combined with position information. Each
decoder block applies normalization, causal self-attention, a residual update, a token-wise
MLP, and another residual update. Causality prevents position $t$ from reading later
tokens. A final vocabulary projection produces one next-token distribution at every
position. For batch size $B$, sequence length $T$, width $D$, and $L$ layers, dense block
projections scale approximately as $O(LBTD^2)$, while pairwise attention scales as
$O(LBT^2D)$. Training memory includes parameters, gradients, optimizer states, and saved
activations; straightforward attention additionally materializes a quadratic $T\times T$
score matrix.

During training, a complete sequence is available at once. If the tokens are
\texttt{[BOS, I, like, tea, EOS]}, the model predicts \texttt{I}, \texttt{like},
\texttt{tea}, and \texttt{EOS} in parallel under a causal mask. The ordinary language-
model loss is
\begin{equation}
  \mathcal{L}_{\mathrm{LM}}
  =-\frac{1}{BT}\sum_{b,t}log p_\theta(x_{b,t+1}\mid x_{b,\leq t}).
  \label{eq:compact-next-token-loss}
\end{equation}
The factor $1/T$ normalizes many token predictions; it does not mean the model performs
only one prediction. Parallel teacher forcing is precisely why training can learn from all
$T$ positions in one forward pass.

Inference executes the same model differently because the next token does not exist until
the preceding token has been sampled. A prefill pass processes the prompt and stores each
layer's keys and values. Each decoding step then computes only the new token's query, key,
and value and attends to the cached prefix. Without this key--value cache, generating token
$4097$ after a 4096-token prompt would recompute the unchanged first 4096 positions; the
same duplication would grow at every later step. The cache removes that repeated prefix
projection at the cost of storing approximately $2LTD$ key and value elements per active
sequence. This distinction matters in post-training because optimization uses parallel
teacher-forced passes, while rollout generation is sequential and often dominates RL
wall-clock time.
